\documentclass{jarticle}
\usepackage[dvipdfmx]{graphicx}
\usepackage{url}
\topmargin       -0.4mm
\oddsidemargin   -0.4mm
\evensidemargin  -0.4mm
\headheight       0.0mm
\headsep          0.0mm
\footskip        10.0mm
%
\textheight     247.0mm
\textwidth      170.0mm
%
\setlength{\parindent}{1zw}
%
\usepackage{multirow}
\usepackage{wrapfig}
\makeatletter
\def \fgcaption{\def \@captype{figure} \caption}
\def \tbcaption{\def \@captype{table} \caption}
\makeatother
\def \figurename{Fig.}
\def \tablename{Table}
%
\def\topfraction{.95}
\def\bottomfraction{.95}
\def\textfraction{.05}
\def\floatpagefraction{.8}
\def\dbltopfraction{.95}
\def\dbltextfraction{.05}
\def\dblfloatpagefraction{.8}
%
\pagestyle{empty}
%
\begin{document}
%
\begin{center}
 \begin{tabular}{|c|c|c|} \hline
  \parbox{50mm}{\vspace{1mm} \centering{\Large\bf 2026/05/28}} & 
  \parbox{55mm}{\vspace{1mm} \centering{\Large\bf B4　堀田　裕樹}} & 
  \parbox{50mm}{\vspace{1mm} \centering{\Large\bf No.~7}}
  \\ [2mm] \hline
%
  \multicolumn{3}{|c|}{%
   \parbox{155mm}{\centering\Large\bf
    \vspace{2mm}
    周期的粗面における超音波反射波の解析とデータ取得
   }
  }
  \\ [3mm] \hline\hline
%
  \multicolumn{3}{|c|}{%
   \parbox{155mm}{\large\vspace{2mm}
%
    {\Large\bf 今月の目標}
%
    \begin{itemize}
     \item 複数パターンの周期粗面試験体に対する実験データの取得と解析
    \end{itemize}
%
    {\Large\bf 進捗の要点}
%
    \begin{itemize}
     \item 試験体 patB および patC の実験データを取得
     \item 時間軸波形の平均化，および各波形の FFT 変換・平均スペクトルの算出を完了
    \end{itemize}
%
    {\Large\bf 今後の実施事項}
%
%
   }
  }
  \\ [3mm] \hline
 \end{tabular}
%
\end{center}
%
\section{はじめに}
今週は，試験体 patB および patC を用いたパルスエコー実験を行い，反射波形データを取得した．
また，既存の patA を含む計3パターンのデータに対して，時間軸および周波数軸での解析を行った．

\section{実験データの解析手順}
取得した各 CSV データに対し，以下の処理を行った．
\begin{enumerate}
    \item \textbf{前処理}: インデックス 0 から 6000 までの範囲（約 3 $\mu$s）を抽出した．
    \item \textbf{時間軸プロット}: 全データの平均波形を算出し，「インデックス vs 電圧」の形式でプロットした．
    \item \textbf{周波数軸解析}: インデックス 6000 までのデータに対し，$2^{14} = 16384$ 点に達するまで末尾をゼロパディングし，FFT 変換を実施した．
\end{enumerate}

\section{解析結果}
Fig. 1 に patA, patB, patC の平均波形を示す．また，Fig. 2 に各試験体における個別の FFT スペクトル（薄線）と平均スペクトル（太線）を示す．

\begin{figure}[htbp]
  \begin{center}
    \begin{tabular}{c}
      \begin{minipage}{0.33\hsize}
        \begin{center}
          \includegraphics[width=\textwidth]{../../tmp_output/exp_waveform_patA.png}
          \\ (a) patA
        \end{center}
      \end{minipage}
      \begin{minipage}{0.33\hsize}
        \begin{center}
          \includegraphics[width=\textwidth]{../../tmp_output/exp_waveform_patB.png}
          \\ (b) patB
        \end{center}
      \end{minipage}
      \begin{minipage}{0.33\hsize}
        \begin{center}
          \includegraphics[width=\textwidth]{../../tmp_output/exp_waveform_patC.png}
          \\ (c) patC
        \end{center}
      \end{minipage}
    \end{tabular}
    \vspace{-3mm}
    \caption{Experimental average waveforms (up to index 6000).}
    \label{fig:waveforms}
  \end{center}
\end{figure}

\begin{figure}[htbp]
  \begin{center}
    \begin{tabular}{c}
      \begin{minipage}{0.33\hsize}
        \begin{center}
          \includegraphics[width=\textwidth]{../../tmp_output/exp_fft_individual_patA.png}
          \\ (a) patA
        \end{center}
      \end{minipage}
      \begin{minipage}{0.33\hsize}
        \begin{center}
          \includegraphics[width=\textwidth]{../../tmp_output/exp_fft_individual_patB.png}
          \\ (b) patB
        \end{center}
      \end{minipage}
      \begin{minipage}{0.33\hsize}
        \begin{center}
          \includegraphics[width=\textwidth]{../../tmp_output/exp_fft_individual_patC.png}
          \\ (c) patC
        \end{center}
      \end{minipage}
    \end{tabular}
    \vspace{-3mm}
    \caption{Individual and average FFT spectra (2--8 MHz).}
    \label{fig:fft_spectra}
  \end{center}
\end{figure}

%参考文献
\begin{thebibliography}{9}
\bibitem{ref1} \url{https://qiita.com/dokagui_tairan/items/11c4451e0283bbb797d8}
\end{thebibliography}

\end{document}
